\newtoggle{withbonus} \toggletrue{withbonus} % set true to show bonus points in grade table

% setting underlying course name (Grundlagenfach Mathematik, §-Unterricht, \ldots)
\newcommand{\mycourse}{Grundlagenfach Informatik}
% name of the class
\newcommand{\myclass}{2. Klassen}
% set date
\newcommand{\mydate}{09. November 2023}

\title{\textbf{Aufgabensammlung}}
\author{Autoren: Patrick Hnilicka, Michael Anderegg, Cyril Wendl, Thomas Graf\\
\mycourse, \myclass}
\date{\mydate}
\maketitle
\thispagestyle{headandfoot}		

\begin{center}
	\fbox{\parbox{140mm}{\centering
			\begin{itemize}
				\item Dauer der Prüfung: $45$ Minuten.
				%\item \textbf{Zeige in jeder Aufgabe unbedingt Deine Schritte!}
				\item Erlaubte Hilfsmittel:
				\begin{itemize}
				    \item Schreibzeug
				    %\item Taschenrechner
				    \item Formelsammlung (Adrian Wetzel)
				\end{itemize}
	\end{itemize}}}
\end{center}
\vspace{10mm}

\begin{center}
	Vorname: \rule{45mm}{0.8pt}\qquad Nachname: \rule{45mm}{0.8pt}
\end{center}
\vspace{20mm}

\begin{center}
	\iftoggle{withbonus} {
		\combinedgradetable[h][questions]
	} {
		\gradetable[h][questions]
	}
\end{center}
\vspace{15mm}
\begin{center}
	Note: 
	\vspace{10mm}
	
	\rule{8mm}{1.2pt}
\end{center}
%\thispagestyle{empty}
\clearpage

\begin{questions}
\section*{Zahlensysteme}
\question[2]{\textbf{Umrechnen in das 5-adische Zahlensystem}}

Wie lautet die kürzeste Darstellung der natürlichen Zahl 72 im 5-adischen System (Basis 5)?
\begin{solution}
Wir berechnen einige Potenzen der Basis 5:
\begin{align*}
    5^3 = 125, \quad 5^2 = 25, \quad 5^1 = 5, \quad 5^0 = 1
\end{align*}
und somit (Ausschöpfen von 72 durch stets möglichst viele der grösstmöglichen Potenzen von 5):
\begin{align*}
    72 = 2\cdot 5^2 + 4\cdot 5^1 + 2\cdot 5^0 = 242_5.
\end{align*}
\end{solution}


\question{\textbf{Zahlensysteme umrechnen}}
\begin{parts}
    \part[1] Wie lautet die kürzeste Darstellung der natürlichen Zahl 53 im 3-adische Zahlensystem (Basis 3)?
    \begin{solution}
        TODO
    \end{solution}
    \part[2] Begründen Sie, weshalb die Ansammlung
    \begin{center}
        \coin{1}\coin{7}\coin{40}\coin{49}
    \end{center}
    von Münzen keine gültige \textit{Münzendarstellung} im 7-adischen Zahlensystem ist.
    \begin{solution}
        TODO
    \end{solution}
    \part[1] Stellen Sie die beiden natürlichen Zahlen 7 und 11 binär dar. Führen Sie danach die binäre schriftliche Addition dieser beiden gefundenen Darstellungen durch. Zeigen Sie Ihre Schritte.
    \begin{solution}
    Die binäre Darstellung von 7 ist $111_2$ und die von 11 ist $1011_2$. Ihre (binäre) schriftliche Addition lautet:
        \begin{center}
            \begin{tabular}{cccccc}
                    & 0 & 0 & 1 & 1 & 1 \\
                +   & 0 & 1 & 0 & 1 & 1 \\
                \hline
                    & 1 & 0 & 0 & 1 & 0
            \end{tabular}
        \end{center}
    \end{solution}
\end{parts}

\section*{Geheimschriften}
\question{\textbf{Zahlensysteme umrechnen}}
\begin{parts}
    \part[1] Ein Text wurde auf zwei unterschiedliche Arten verschlüsselt: einmal mit der Methode Caesar und einmal mit der Methode Vigenère. Betrachten Sie die folgenden beiden Darstellungen der Friedmannsch'en Charakteristik für unterschiedliche Schlüssellängen und bestimmen Sie, zu welcher Verschlüsselung sie gehören:
    ...TODO, add figures of FC for Caesar and Vigenere.\\
    
    \begin{solution}
        TODO
    \end{solution}
    
\end{parts}
\question Was berechnet die Friedmann'sche Charakteristik? Kreuzen Sie die richtige Antwort an:\\
\begin{oneparchoices}
 \choice Die Friedmann'sche Charakteristik $FC$ eines Texts $T$ ($FC_T$) berechnet die Schlüssellänge durch als die Summe der Abweichungen (im Quadrat) der Häufigkeit aller Zahlen von der erwarteten Häufigkeit (1/26). Für deutsche Text ist $FC_T\sim 0.0388$. \\
 \choice Die Friedmann'sche Charakteristik $FC$ eines Texts $T$ ($FC_T$) berechnet die Schlüssellänge, indem die häufigsten Bi- oder Trigramme in einem Text gesucht werden und deren Abstände berechnet und in Primfaktoren zerlegt werden. Die am häufigsten vorkommenden Primfaktoren bestimmen die Schlüssellänge.\\
 \choice Die Friedmann'sche Charakteristik erlaubt es uns, mit Vigenère verschlüsselte Geheimtexte zurück in den Klartext umzuwandeln, indem eine Häufigkeitsanalyse im ursprünglichen Text durchgeführt wird. Der am häufigsten vorkommende Buchstabe im Geheimtext entspricht bei genug langer Textlänge erwartungsgemäss dem häufigsten Buchstaben in der jeweiligen Sprache des Klartext.
 \begin{solution}
     Die Friedmann'sche Charakteristik $FC$ eines Texts $T$ ($FC_T$) berechnet die Schlüssellänge durch als die Summe der Abweichungen (im Quadrat) der Häufigkeit aller Zahlen von der erwarteten Häufigkeit (1/26). Für deutsche Texte ist $FC_T\sim 0.0388$.
 \end{solution}
\end{oneparchoices}
\droptotalpoints
\end{questions}
